\section{Derivative prices for the Gaussian correction chains}
\label{sec:itinerary-prices}
The corrections in Section~\ref{sec:local-covariance} are sums
over ordered itineraries of physical bonds. Their operator bounds alone do not price each
activity in the connected expansion. We now keep the itinerary
labels and their full read supports while differentiating and
summing them.

\subsection{Products along a fixed itinerary}
Write $q$ for the number of colour coordinates. An itinerary of
length $N$ consists of physical bonds $e_0,\ldots,e_N$; colour
indices are summed within its matrix product. Suppose the entry
jets of its $r$th matrix have amplitude $a_r\ge0$ and common
derivative scale $d\ge0$. Repeated matrix multiplication and the
product rule give
\begin{equation}\label{eq:itinerary-entry-price}
 \mathcal B_j\!\left(q^N\prod_{r=0}^{N-1}a_r,Nd\right)
\end{equation}
for any entry of the colour-block product. This follows by
induction on $N$: multiplying the next entry adds $d$ to the
derivative scale, multiplies the amplitude by $a_r$, and introduces
at most $q$ choices for the next colour. The estimate deliberately
overcounts when $N=1$. Taking a colour trace adds a factor $q$.
Integrating such a trace over $s\in[0,1]$ preserves this bound
when the entry estimates are uniform in $s$.

For the present matrices, this last differentiation under the
integral is valid by the compact-glide argument of
Section~\ref{sec:normalization-jets}: all principal determinants
are nonzero on one common parameter neighbourhood for $s\in[0,1]$,
and the finite matrix products are jointly smooth there. A closed
itinerary contributes its diagonal trace; a nonclosed itinerary
contributes zero to that trace.

All the matrices in these words have range at most $3B$. If some
consecutive pair of bonds has all endpoint distances greater than
$3B$, the corresponding matrix block vanishes identically as the
coarse field varies. Thus the whole word is a zero germ. Call an
itinerary admissible when every consecutive pair has endpoints
at distance at most $3B$. This condition is fixed by its geometry,
and is retained in every price below.

Let $a$ be the common weighted jet amplitude of the actual
precision $A$, let $h_0$ price the padded inverse of the glide,
and let $t_0$ price its error $A_sP_s-1$. These quantities were
constructed in Section~\ref{sec:normalization-jets}. The amplitudes
$a,h_0$ have upper exponent zero; $t_0$ has every negative upper
exponent. The first logarithmic correction has word
\[
 (A-1)P_s(-E_s)^{n+1},\qquad E_s=A_sP_s-1.
\]
It has $n+3$ matrix factors. If $d_A$ bounds their derivative
scales, \eqref{eq:itinerary-entry-price} and the colour trace give
\begin{equation}\label{eq:itinerary-A-price}
 \mathcal B_j\bigl(C_A Q_A^n,(n+3)d_A\bigr),
 \qquad C_A=q^4(a+1)h_0t_0,\quad Q_A=qt_0.
\end{equation}
The same argument applies to the flat-reference correction.
Its coarse field is constant and satisfies the same entry bounds.

For the second correction put $D=AC-1$, where $C$ is the symmetric
padded covariance. Its entry jets have amplitude $w_0$ with every
negative upper exponent, obtained by multiplying the padded-inverse
error by the full inverse envelope and then by $A$. If $d_D$ is
their derivative scale, the word $D(-sD)^n$ has $n+1$ factors.
Since $0\le s\le1$, its integrated trace has bound
\begin{equation}\label{eq:itinerary-S-price}
 \mathcal B_j\bigl(C_S Q_S^n,(n+1)d_D\bigr),
 \qquad C_S=q^2w_0,\quad Q_S=qw_0.
\end{equation}
These formulas price the actual trace words, including their
endpoint identifications. Discarding a factor $1/2$ or another
coefficient of absolute value at most one only enlarges the bounds.

\subsection{The two protected Gaussian endpoints}
The inverse correction is a quadratic form with word
$(-D)^{n+1}A$, of length $n+2$, and a Gaussian coordinate at
each endpoint. Those coordinates are not bounded on the whole
sample space. Their protective cutoffs supply the needed local
bound.

Let $u(y)=\nu(y)+g_\beta z(y)$, $g_\beta=\beta^{-1/2}$,
be the compensated increment on an active physical bond.
If its doubled cutoff is not a zero germ, the increment jet bound
has amplitude $s_0+4a_0\tau$, where $s_0$ has upper exponent
$-1/2$ and $a_0>0$ is the fixed cutoff parameter. The mean jets
have amplitude $m_0$ of upper exponent $-1/2$. Hence
\begin{equation}\label{eq:itinerary-endpoint-price}
 z_e\in\mathcal B_j(p,1),\qquad
 p=\sqrt\beta\,(s_0+4a_0\tau+m_0),
\end{equation}
which has upper exponent zero since $\beta\le2h$.
If either cutoff is a zero germ, the protected quadratic activity
is a zero germ instead. If an endpoint is inactive, the inactive
block of the completed precision and covariance is the identity;
the corresponding defect block is zero. These block identities
persist on the neighbourhood where the fixed principal windows
remain invertible. The quadratic word therefore vanishes as a
germ unless both endpoints are active. This justifies using
\eqref{eq:itinerary-endpoint-price} only where it applies.

Suppose the cutoff jets have amplitude one and scale $d_0$.
The word has entry amplitude $q^{n+2}aw_0^{n+1}$.
The two endpoint colour sums add $q^2p^2$, and the product
rule adds two endpoint scales and two cutoff scales. Enlarging
the resulting scale gives
\begin{equation}\label{eq:itinerary-quadratic-price}
 \mathcal B_j\bigl(C_Q Q_Q^n,(n+2)d_Q\bigr),\qquad
 C_Q=p^2q^4aw_0,\quad Q_Q=qw_0,\quad
 d_Q=d_D+d_0+1,
\end{equation}
after enlarging $d_D$ to include the scale of $A$.
All derivative scales here have upper exponent zero. The
unprotected linear score is handled separately: its coefficient
has every negative upper exponent by the mean equation, while
its one protected endpoint has the bound just proved. Summing
its $q$ colours and multiplying its cutoff produces one local
bond activity with every negative upper exponent, even after
its fixed-radius support is counted.

\subsection{Paying the derivative order before the chain length}
All three itinerary families have the form
\[
 \mathcal B_j(CQ^n,(n+u)d),\qquad u\in\{3,1,2\},
\]
on admissible itineraries and zero germs otherwise. To give every
jet a common amplitude and unit derivative scale, multiply the
amplitude by $(1+(n+u)d)^j$. The elementary inequality
$n+1\le2^n$ implies
\[
 1+(n+u)d\le2^n\{1+(1+u)d\}.
\]
It is therefore enough to replace
\begin{equation}\label{eq:itinerary-inflated-price}
 C\longmapsto C'=C\{1+(1+u)d\}^{j},\qquad
 Q\longmapsto Q'=2^jQ.
\end{equation}
No factor depending on $n$ remains outside the geometric power.
For fixed $j$, $C'$ retains every negative upper exponent of $C$,
and $Q'$ retains every negative upper exponent of $Q$.

\subsection{The rooted support sum}
Fix a halo multiplier $m\ge4$, also large enough to contain the
previously established mean, window and compensated-variation read
radii, a support exponent $a_1\ge0$, and a length reserve
$\varepsilon>0$. The inequality $m\ge4$ alone is used in the
count below; the additional read-radius requirement is needed
when applying it to the actual activities. The support of an itinerary
is the union of the radius-$mB$ halos of the endpoint sets of
all its bonds. Put
\[
 \rho_L=L^4\#\{\mathrm{axes}\},\quad
 K=b_{3B+2}\rho_L,\quad s_m=b_{m+1}.
\]
There are at most $K$ choices for a next bond in an admissible
step: its anchor is within $3B+2$ of the preceding anchor.
Thus there are at most $K^{n+u}$ admissible itineraries of length
$n+u$ with prescribed initial bond.

Each bond halo has $B$-cover number at most $s_m$. Consecutive
halos intersect because $3B+2\le mB$. Subadditivity along the
ordered chain gives support cover at most $(n+u+1)s_m$.
If this support contains a fixed site $x_0$, the initial anchor
is within $(n+u+1)(mB+1)$ of $x_0$: follow the chain to a halo
containing $x_0$ and sum the step distances. Since
$n+u+1\le(u+1)2^n$, the number of possible initial bonds is at
most
\[
 b_{(u+1)(mB+1)}16^n\rho_L.
\]
The factor $16^n$ is the fourth power of $2^n$; it also bounds
the additive one in the lattice ball formula. Wrapping on a torus
can only reduce these counts.

Multiplying the price $C'(Q')^n$, the count, the support weight
and the extra length weight $e^{2\varepsilon n}$ gives the bound
\begin{equation}\label{eq:itinerary-site-series}
 \sum_{\lambda:x_0\in S_\lambda}
 p_\lambda e^{a_1\operatorname{cov}_B(S_\lambda)}
                   e^{2\varepsilon n(\lambda)}
 \le H C'\sum_{n\ge0}
 \bigl[Q'16K e^{a_1s_m+2\varepsilon}\bigr]^n,
\end{equation}
where
\[
 H=b_{(u+1)(mB+1)}\rho_L K^u e^{a_1(u+1)s_m}.
\]
All sums have nonnegative terms, so grouping first by length and
then by initial bond is justified without a convergence assumption.
At fixed $m,L,j$, both $K$ and $H$ have upper exponent zero.
The geometric ratio in \eqref{eq:itinerary-site-series} is
eventually at most $1/2$, and its sum is at most $2HC'$.
Consequently every prescribed negative upper exponent is available
for the rooted price of these correction chains. This pays the
actual full read supports and every itinerary label, including
repeated physical bonds; it does not assume disjoint paths or
independent Gaussian coordinates.
